% ===================================================================
%  Dodatek L: Formuła sumy krytycznych sprzężeń kinetycznych
%             i predykcja stosunku mas leptonów
%  TGP v1 — Sesja v34 (2026-03-28)
%  Ścieżka 10: skan alpha_kin, amplituda ogona, sum formula
% ===================================================================
\section{Formuła sumy: skan kinetyczny i~stosunek mas leptonów}%
\label{app:formula-sumy}

Niniejszy dodatek dokumentuje \textbf{Ścieżkę~10} (skrypty
\texttt{ex62}--\texttt{ex85}): podejście, w~którym kinetyczny parametr
sprzężenia $\alpha_{\rm kin}$ jest skanowany, a~warunek
\begin{equation}\label{eq:L-def-warunek}
  \left(\frac{A_{\rm tail}(\varphi \cdot z_0(\alpha_{\rm kin}))}
              {A_{\rm tail}(z_0(\alpha_{\rm kin}))}\right)^{\!4}
  = r_{21} = \frac{m_\mu}{m_e} \approx 206{,}768
\end{equation}
wyznacza \textbf{dwie krytyczne wartości} $\alpha^*_1$, $\alpha^*_2$,
których suma jest zbliżona do prostej kombinacji stałych fundamentalnych TGP:
\begin{equation}\label{eq:L-sum-formula}
  \boxed{\alpha^*_1 + \alpha^*_2
  \;\approx\; 2\pi - \frac{11}{10}
  \;=\; 2\pi - \frac{\alpha_{\rm TGP}}{2} - \frac{\beta_{\rm pot}}{10}.}
\end{equation}
gdzie $\alpha_{\rm TGP} = 2$ (parametr substratowy N0-2) i~$\beta_{\rm pot} = 1$
(współczynnik potencjału w~jednostkach $\beta = \gamma = 1$, N0-5).

\begin{remark}[Duch teorii: brak nowych parametrów]
Skan $\alpha_{\rm kin}$ \textbf{nie wprowadza} nowego wolnego parametru.
Sprzężenie $\alpha_{\rm TGP} = 2$ pozostaje ustalone. Skan jest
\emph{narzędziem diagnostycznym}: pyta, przy jakich wartościach efektywnego
kinetycznego sprzężenia warunek~\eqref{eq:L-def-warunek} jest spełniony.
Wynikające $\alpha^*_{1,2}$ nie są parametrami teorii, lecz cechami spektrum
operatora solitonowego TGP przy danym $r_{21}$. Formuła~\eqref{eq:L-sum-formula}
jest hipotezą o~strukturze tego spektrum.
\end{remark}

% -------------------------------------------------------------------
\subsection{ODE solitonu przy zmiennym sprzężeniu kinetycznym}
\label{app:L-ode}
% -------------------------------------------------------------------

\subsubsection{Równanie bazowe}

ODE solitonu sferycznie symetrycznego (por.\ dodatek~\ref{app:ogon-masy}):
\begin{equation}\label{eq:L-ode}
  f_\alpha(g)\,g'' + \frac{2}{r}\,g' = V'(g),
  \qquad
  f_\alpha(g) = 1 + 2\alpha_{\rm kin}\ln g,
  \qquad
  V'(g) = g^2(1-g),
\end{equation}
z~warunkami: $g(0) = g_0$, $g'(0) = 0$, $g(r)\to 1$ dla $r\to\infty$.

\textbf{Uwaga fizyczna:} dla $\alpha_{\rm kin} = 2$ (wartość TGP)
otrzymujemy identyczne ODE jak w~dodatkach~J--K. Skan po $\alpha_{\rm kin}$
bada \emph{rodzinę} ODE sparametryzowaną tym współczynnikiem.

\subsubsection{Punkt osobliwy i~bariera kinetyczna}

Czynnik $f_\alpha$ zeruje się przy:
\begin{equation}\label{eq:L-gstar}
  g^*(\alpha_{\rm kin}) = \exp\!\left(-\frac{1}{2\alpha_{\rm kin}}\right),
  \qquad
  g^*(2) = e^{-1/4} \approx 0{,}779.
\end{equation}
Bariera kinetyczna ($g \geq g^*$) obowiązuje dla każdego $\alpha_{\rm kin} > 0$.

\subsubsection{Punkt strzału $z_0(\alpha_{\rm kin})$: czysta sinusowa faza}

Spośród ciągłego spektrum $g_0 \in (g^*, \infty)$ wyróżniamy punkt
\textbf{czystej fazy sinusowej}:
\begin{equation}\label{eq:L-z0-def}
  z_0(\alpha_{\rm kin})\colon\qquad B_{\rm tail}(z_0, \alpha_{\rm kin}) = 0,
\end{equation}
gdzie $B_{\rm tail}$ jest amplitudą kosinusową ogona (por.\ dodatek~J,
hipoteza H1, stwierdzenie~\ref{prop:J-quantization}).
$z_0$ jest pierwszym zerem $B_{\rm tail}$ w~sektorze $n_{\rm bounce}=0$.

\begin{remark}[Wartości numeryczne]
Dla fizycznego $\alpha_{\rm kin} = 2$:
$z_0(2) = g_0^e \approx 1{,}230$ (elektron, wyniki \texttt{ex58}, 0{,}8\% od $1{,}24$).
\end{remark}

% -------------------------------------------------------------------
\subsection{Warunek krytyczny i~dwie gałęzie}
\label{app:L-warunek}
% -------------------------------------------------------------------

\subsubsection{Funkcja testowa}

Definiujemy:
\begin{equation}\label{eq:L-F-def}
  F(\alpha_{\rm kin}) \;\equiv\;
  \left(\frac{A_{\rm tail}(\varphi\cdot z_0(\alpha_{\rm kin}),\,\alpha_{\rm kin})}
              {A_{\rm tail}(z_0(\alpha_{\rm kin}),\,\alpha_{\rm kin})}\right)^{\!4},
\end{equation}
gdzie $\varphi = (1+\sqrt{5})/2 \approx 1{,}618$ jest złotą proporcją,
a~amplituda $A_{\rm tail}$ jest wyznaczana przez ekstrapolację z~okien
$r \in [r_L, r_L + \Delta r]$ (metoda \texttt{curve\_fit} w~\texttt{ex62}).

\begin{remark}[Rola złotej proporcji]
Czynnik $\varphi$ pochodzi z~obserwacji $g_0^\mu / g_0^e \approx \varphi$
(hipoteza~\ref{hyp:J-golden-topo}). Wartość $\varphi\cdot z_0$ odpowiada
wówczas \emph{skalowaniu muonowemu} punktu kwantowania elektronu.
\end{remark}

\subsubsection{Dwie krytyczne wartości $\alpha^*$}

\begin{result}[Krytyczne wartości $\alpha^*$ --- wyniki finalne (\texttt{ex88}, v36)]%
\label{res:L-alpha-stars}
Numeryczne poszukiwanie zer $F(\alpha_{\rm kin}) - r_{21} = 0$
w~przedziale $\alpha_{\rm kin} \in [2{,}0;\, 3{,}5]$
ze~sta\l{}ym oknem $B_{\rm coeff}\in[28,42]$
(skrypt \texttt{ex88}, $r_{\rm max}\in\{100,120,150,200,300\}$)
daje \textbf{dok\l{}adnie dwa rozwiązania}, identyczne dla wszystkich
$r_{\rm max}$ (dok\l{}adno\'s\'c maszynowa):
\begin{equation}\label{eq:L-alpha-stars}
  \alpha^*_{1,\infty} = 2{,}43183767\,,
  \qquad
  \alpha^*_{2,\infty} = 2{,}63557742\,.
\end{equation}

\textbf{Historia i~korekty numeryczne}:
\begin{itemize}[nosep]
  \item \texttt{ex79}--\texttt{ex80}: $\alpha^*_1 \approx 2{,}4414$, suma $5{,}1832$,
    odchylenie $0{,}13\,\text{ppm}$ --- wynik z~artefaktu grubej siatki (ex71--ex83).
  \item \texttt{ex84} ($r_{\rm max}=120$, sta\l{}e okno): $\alpha^*_1 \approx 2{,}440$,
    suma $5{,}1349$, odchylenie $9307\,\text{ppm}$ --- pierwsze prawid\l{}owe.
  \item \texttt{ex87v2} ($r_{\rm max}=100$, okno [28,42]): $\alpha^*_1=2{,}4318$,
    $\alpha^*_2=2{,}6356$, suma $5{,}0674$, $22336\,\text{ppm}$.
  \item \texttt{ex88} (sta\l{}e okno, 5 warto\'sci $r_{\rm max}$): \textbf{potwierdzenie}
    $22336\,\text{ppm}$; $\sigma_{\alpha^*} = 0$ --- ca\l{}kowita stabilno\'s\'c.
\end{itemize}

\textbf{Wcze\'sniejsze $\alpha^* \approx 2{,}44$ i~$2{,}74$} (ex71--ex80) by\l{}y
artefaktem skalowanego okna $B_{\rm coeff}$ (O-L1 rozwi\k{a}zane).
\end{result}

% -------------------------------------------------------------------
\subsection{Formuła sumy}
\label{app:L-suma}
% -------------------------------------------------------------------

\begin{hypothesis}[Formuła sumy krytycznych sprzężeń]%
\label{hyp:L-sum-formula}
Suma $\alpha^*_1 + \alpha^*_2$ jest dokładnie wyrażona przez
substratowe stałe TGP:
\begin{equation}\label{eq:L-sum-exact}
  \alpha^*_1 + \alpha^*_2
  \;=\; 2\pi - \frac{\alpha_{\rm TGP}}{2} - \frac{\beta_{\rm pot}}{10}
  \;=\; 2\pi - 1 - \frac{1}{10}
  \;=\; 2\pi - \frac{11}{10},
\end{equation}
gdzie:
\begin{itemize}[nosep]
  \item $\alpha_{\rm TGP} = 2$ (substratowy wykładnik kinetyczny, N0-2),
  \item $\beta_{\rm pot} = 1$ (normalizacja potencjału $V'(g) = \beta_{\rm pot}\,g^2(1-g)$, N0-5).
\end{itemize}
Numerycznie: $2\pi - 11/10 = 5{,}18318\ldots$

\textbf{Motywacja}: skrypt \texttt{ex73} weryfikował ogólność formuły
$S = 2\pi - 1 - \beta_{\rm pot}/10$ dla różnych wartości $\beta_{\rm pot}$.
Jeśli formuła zachodzi dla każdego $\beta_{\rm pot}$, to jest ona
\textbf{prawem strukturalnym} --- nie numerologią dla $\beta_{\rm pot} = 1$.

Status: \statuslabel{Hipoteza} \textbf{obalona} (v36, O-L1 zamkni\k{e}te):
$S(\infty) = 5{,}0674 \neq 2\pi{-}11/10 = 5{,}1832$ ($22336\,\text{ppm}$,
\texttt{ex88}).
Brak prostej formu\l{}y algebraicznej dla~$S$.
Otwarte: (ii)~weryfikacja ogólno\'sci przez skan $\beta_{\rm pot}$
(\texttt{ex89}, O-L2), (iii)~analityczna derywacja $\alpha^*_{1,2}$.
\end{hypothesis}

\begin{remark}[Związek z~$\alpha_{\rm TGP}$ i~$2\pi$]
Formuła~\eqref{eq:L-sum-exact} zawiera $2\pi$ --- pełny cykl fazy
w~opisie solitonu oscylacyjnym. Człon $\alpha_{\rm TGP}/2 = 1$ jest
\emph{połową} substratowego wykładnika kinetycznego; człon $1/10$ pochodzi
z~normalizacji potencjału próżniowego.

Fizyczna interpretacja: $\alpha^*_1$ i~$\alpha^*_2$ dzielą okrąg $2\pi$
tak, że różnica ich odległości od $\pi$ (centrum) jest wyznaczona przez
stałe substratowe $\alpha_{\rm TGP}$ i~$\beta_{\rm pot}$.
Jest to \textbf{warunek półsumy}: struktura spektrum $F(\alpha)$ jest
symetrycznie eksponowana wokół $\pi - (11/10)/2 = \pi - 11/20$.
\end{remark}

% -------------------------------------------------------------------
\subsection{Analiza algebraiczna}
\label{app:L-algebra}
% -------------------------------------------------------------------

\begin{result}[Własności algebraiczne $\alpha^*_1$, $\alpha^*_2$ z~\texttt{ex80}]%
\label{res:L-algebra}
Przy wartościach z~\texttt{ex79}: $\alpha^*_1 = 2{,}44143$,
$\alpha^*_2 = 2{,}74175$, minimum $\alpha_{\min} \approx 2{,}5633$,
$F_{\min} \approx 197{,}95$:
\begin{center}\footnotesize
\begin{tabular}{@{}lll@{}}
\toprule
\textbf{Wielkość} & \textbf{Wartość numeryczna} & \textbf{Formuła kandydacka} \\
\midrule
$S = \alpha^*_1 + \alpha^*_2$ & $5{,}18318$ & $2\pi - 11/10$ (0{,}13 ppm) \\
$D = \alpha^*_2 - \alpha^*_1$ & $0{,}30032$ & $3/10$ (0{,}1\%) \\
$P = \alpha^*_1 \cdot \alpha^*_2$ & $6{,}7004$ & $(\pi - 7/10)(\pi - 4/10)$ (0{,}01\%) \\
$\alpha_{\min}$    & $2{,}5633$  & $(S-D)/2 = S/2 - 3/20$ \\
$F_{\min}$         & $197{,}95$  & brak prostej formuły \\
\bottomrule
\end{tabular}
\end{center}

\textbf{Uwaga}: iloczyn $P \approx (\pi-7/10)(\pi-4/10)$ sugeruje,
że $\alpha^*_1 = \pi - 7/10 + \epsilon_1$ i~$\alpha^*_2 = \pi - 4/10 + \epsilon_2$
z~małymi korektami $\epsilon_{1,2}$, których suma $\epsilon_1 + \epsilon_2$
pochodzi od członu $-(11/10 - (7/10 + 4/10)) = 0$ (dokładnie),
a~więc formuła sumy $S = 2\pi - 11/10$ wynika z~centrosymetrii
$\alpha^*_{1,2} = \pi - (11/20) \pm \Delta$ z~$\Delta = D/2 = 3/20$.

Status: \statuslabel{Hipoteza} algebraiczna (\texttt{ex80}: 3/5 PASS dla prostych
formuł; wyniki niefinalnie ze względu na O-L1).
\end{result}

% -------------------------------------------------------------------
\subsection{Ogólność formuły: weryfikacja skanem $\beta_{\rm pot}$}
\label{app:L-generalnosc}
% -------------------------------------------------------------------

\begin{result}[Skan $\beta_{\rm pot}$ z~\texttt{ex73}]%
\label{res:L-beta-scan}
Skrypt \texttt{ex73} testował formułę $S(\beta_{\rm pot}) = 2\pi - 1 - \beta_{\rm pot}/10$
dla wartości $\beta_{\rm pot} \in \{0{,}5;\, 1{,}0;\, 1{,}5;\, 2{,}0\}$
(gdzie $\beta_{\rm pot} = 1$ jest wartością fizyczną TGP).

Wynik (skrypt \texttt{ex73}, metoda pełna):

\begin{center}\footnotesize
\begin{tabular}{@{}ccccc@{}}
\toprule
$\beta_{\rm pot}$ & $\alpha^*_1$ & $\alpha^*_2$ & $S$ numeryczna & $2\pi - 1 - \beta/10$ \\
\midrule
$0{,}5$ & do wyznaczenia & do wyznaczenia & do wyznaczenia & $5{,}233$ \\
$1{,}0$ & $2{,}4414$ & $2{,}7418$ & $5{,}1832$ & $5{,}183$ \checkmark \\
$1{,}5$ & do wyznaczenia & do wyznaczenia & do wyznaczenia & $5{,}133$ \\
$2{,}0$ & do wyznaczenia & do wyznaczenia & do wyznaczenia & $5{,}083$ \\
\bottomrule
\end{tabular}
\end{center}

Wyniki dla $\beta_{\rm pot} \neq 1$ zależą od zbieżności solvera;
pełne wyniki wymagają uruchomienia \texttt{ex73} na docelowym systemie.

Status: \statuslabel{Program} (weryfikacja ogólności; O-L2).
\end{result}

% -------------------------------------------------------------------
\subsection{Interpretacja fizyczna}
\label{app:L-fiz}
% -------------------------------------------------------------------

\begin{remark}[Ścieżka 10 vs Ścieżka 9]
Ścieżka~9 (Dod.~J--K) identyfikuje masy leptonów z~amplitudą ogona
przy \textbf{stałym $\alpha_{\rm kin} = \alpha_{\rm TGP} = 2$} i~\emph{różnych}
$g_0^{e,\mu,\tau}$ (warunki początkowe ODE).

Ścieżka~10 (niniejszy dodatek) identyfikuje masy leptonów poprzez
\textbf{skan $\alpha_{\rm kin}$} przy \emph{stałym} $g_0/z_0 = \varphi$ (złota proporcja).

Obie ścieżki są komplementarne: Ścieżka~9 zakłada, że lepton wybiera
$g_0$, a~Ścieżka~10 zakłada, że lepton odpowiada specjalnemu efektywnemu
sprzężeniu kinetycznemu. Fizyczna unifikacja obu podejść --- tzn.\
dowód, że $g_0$ wyznaczony z~Ścieżki~9 jest spójny z~$\alpha^*$ z~Ścieżki~10
--- jest programem obliczeniowym (O-L3).
\end{remark}

\begin{hypothesis}[Interpretacja $\alpha^*$ jako efektywnego sprzężenia]%
\label{hyp:L-alpha-eff}
Wartości $\alpha^*_1$ i~$\alpha^*_2$ mogą odpowiadać
\textbf{efektywnym stałym sprzężenia kinetycznego} solitonu
w~różnych sektorach topologicznych:
\[
  \alpha^*_1 \approx \alpha_{\rm TGP} + \frac{4\ln(g_0^e)}{2\ln(g_0^e)} \cdot f_1
  \qquad \text{(sektor elektronu)},
\]
\[
  \alpha^*_2 \approx \alpha_{\rm TGP} + \Delta_\mu
  \qquad \text{(sektor mionu)},
\]
gdzie $\Delta_\mu$ pochodzi od korekcji wynikających z~$n_{\rm cross} = 1$
(jeden przeskok przez poziomicę próżni).

\textbf{Konsekwencja}: stosunek mas $r_{21} = (A(\varphi z_0)/A(z_0))^4$
przy \emph{odpowiednio dobranych} $\alpha^*$ jest \textbf{tożsamością
strukturalną} TGP --- wnioskiem z~aksjomat\'ow N0-2 i~N0-5,
a~nie z~wolnego parametru.

Status: \statuslabel{Hipoteza} (wymaga O-L3 i~O-L1).
\end{hypothesis}

% -------------------------------------------------------------------
\subsection{Otwarte problemy Ścieżki~10}
\label{app:L-otwarte}
% -------------------------------------------------------------------

\begin{center}\footnotesize
\begin{tabular}{@{}p{1.0cm}p{4.5cm}lp{4.0cm}@{}}
\toprule
\textbf{ID} & \textbf{Problem} & \textbf{Status} & \textbf{Oczekiwany wynik} \\
\midrule
O-L1 & \textbf{Zamkni\k{e}te} (\texttt{ex88}, v36, 4/4~PASS).
       Okno $B_{\rm coeff} \in [28, 42]$ sta\l{}e $\Rightarrow$
       $\alpha^*_{1,2}$ identyczne dla
       $r_{\rm max}\in\{100,120,150,200,300\}$ (dok\l{}adno\'s\'c maszynowa).
       \textbf{Warto\'sci finalne}:
       $\alpha^*_{1,\infty} = 2{,}43183767$,
       $\alpha^*_{2,\infty} = 2{,}63557742$,
       $S(\infty) = 5{,}06742$.
       Brak prostej formu\l{}y algebraicznej dla~$S$;
       najbli\.zszy kandydat $2\pi - 11/10$ odst\k{a}puje o~$22336\,\text{ppm}$.
     & \statuslabel{Zamkni\k{e}te}
     & Obalenie $S = 2\pi{-}11/10$ potwierdzone definitywnie
       ($22\,336\,\text{ppm}$). \\[4pt]
O-L2 & \textbf{Zamkni\k{e}te} (\texttt{ex89}, v36).
       Skan $\beta_{\rm pot} \in \{0{,}5;\, 0{,}8;\, 1{,}0;\, 1{,}2;\, 1{,}5\}$
       ze~sta\l{}ym oknem $[28,42]$, $r_{\rm max}=120$:
       \textbf{tylko $\beta_{\rm pot}=1{,}0$ daje dwa zera}
       ($\alpha^*_1=2{,}4316$, $\alpha^*_2=2{,}6356$, $S=5{,}0672$).
       Dla $\beta_{\rm pot} \neq 1$: $F(\alpha_{\rm kin}) < R_{21}$
       lub brak zero w~$[2{,}0;\,3{,}5]$ --- zera \textbf{nie s\k{a} topologiczne}.
       Interpretacja: $\beta_{\rm pot}=1$ jest warunkiem pr\'o\.zniowym
       N0-5 ($\beta=\gamma$), wi\k{e}c sektor leptonowy wymaga
       dok\l{}adnie fizycznego potencja\l{}u TGP.
     & \statuslabel{Zamkni\k{e}te}
     & Zera specyficzne dla $\beta_{\rm pot}=1$ (fizycznej warto\'sci TGP).
       hyp:L-alpha-eff: bez awansu. \\[4pt]
O-L3 & \textbf{Zamkni\k{e}te cz\k{e}\'sciowo} (\texttt{ex91}, v36, 4/5~PASS).
       \'Scie\.zka~9 ($\alpha=2$): $z_0(2)=1{,}228$,
       $F(2)=263{,}9 \neq R_{21}=206{,}8$ ($+27{,}7\%$, FAIL).
       \'Scie\.zka~10 ($\alpha^*_1$): $z_0(\alpha^*_1)=1{,}258$,
       $F(\alpha^*_1)=206{,}8$ (PASS).
       $z_0(\alpha^*_1) \neq g_0^e(2)$ (r\'o\.znica $2{,}4\%$);
       $A_{\rm tail}^{P10}/A_{\rm tail}^{P9} = 1{,}20$.
       \textbf{Wniosek}: \'scie\.zki komplementarne, nie identyczne ---
       $\alpha_{\rm TGP}=2$ nie spe\l{}nia warunku $F=R_{21}$
       bezpo\'srednio; potrzebne $\alpha^*_1 \neq \alpha_{\rm TGP}$.
     & \statuslabel{Zamkni\k{e}te}
     & Unifikacja cz\k{e}\'sciowa; pytanie O-L4: sk\k{a}d $\alpha^*\neq\alpha_{\rm TGP}$? \\[4pt]
O-L4 & \textbf{Mechanizm cz\k{e}\'sciowo wyja\'sniony} (\texttt{ex92}, v37).
       Profil $F(\alpha)$ na ga\l{}\k{e}zi fizycznej ($z_0\approx 1{,}22$--$1{,}26$):
       $F(2)=264{,}1 > R_{21}=206{,}77$; minimum $F\approx201{,}0$
       w~$\alpha_{\min}\approx 2{,}49$--$2{,}54$ (midpoint zer: $(2{,}4345+2{,}6441)/2=2{,}539$),
       po\l{}\k{e}te za $\alpha_{\rm TGP}=2$.
       \textbf{Wniosek}: $\alpha_{\rm TGP}=2$ le\.zy \textbf{przed} minimum $F$
       (ga\l{}\k{a}\'z malej\k{a}ca); zera $\alpha^*_{1,2}$ otaczaj\k{a} minimum.
       \textbf{Otwarte}: analityczna formu\l{}a dla~$\alpha_{\min}$
       i~$z_0(\alpha_{\min})$.
     & \statuslabel{Program}
     & Formu\l{}a $\alpha_{\min}\approx2{,}54$; analitycznie: warunek $dF/d\alpha=0$ \\[4pt]
O-L5 & \textbf{$\alpha^*_3$ --- definitywnie brak rozwi\k{a}zania}
       (\texttt{ex95}--\texttt{ex102}, v38--v39).
       $G_3(\alpha)\equiv(A_{\rm tail}(m\cdot z_0)/A_{\rm tail}(z_0))^4$;
       \textbf{Wyniki v38}: bias okien STD wykryty (\texttt{ex95});
       $G_{3,\rm far}(\varphi^2 z_0)>R_{31}$ na $\alpha\in[2{,}5;\,5{,}0]$
       --- min $4667$ przy $\alpha=2{,}50$ ($+34\%$, \texttt{ex96});
       plateau $A_\tau\approx2{,}994$ potwierdzone do $r=286$ (\texttt{ex97}).
       \textbf{Wyniki v39}: skan $\alpha\in[1{,}5;\,2{,}55]$ --- min $G_{3,\rm far}=4511$
       przy $\alpha=2{,}35$ ($+30\%$, \texttt{ex100}).
       Pozorne zero przy $\alpha\approx7{,}1$ to artefakt okna $[60,76]$
       --- prawdziwe $G_3(\alpha{=}7{,}0)\approx65000$ (\texttt{ex101}).
       Test 11 mno\.znik\'ow $m$ z~oknami VFAR $[80,156]$ (\texttt{ex102}):
       najbli\.zej $R_{31}$: $m=2{,}5$ przy $\alpha=\alpha^*_{1,\rm far}$
       daje $G_3=3519$ ($+1{,}2\%$) --- brak zera.
       $G_{3,\rm vfar}(\varphi^2 z_0)>R_{31}$ na $\alpha\in[2{,}3;\,5{,}0]$
       --- min $4519$ przy $\alpha=2{,}30$.
       \textbf{Wniosek definitywny}: formu\l{}a amplitudowa
       $G_3=(A_\tau(m\cdot z_0)/A_e(z_0))^4=R_{31}$ nie ma rozwi\k{a}zania
       dla \.zadnego $m\in\{\varphi,\sqrt{2},\ldots,\varphi^3\}$.
       Sektor $\tau$ wymaga nowej parametryzacji.
     & \statuslabel{Program}
     & $\alpha^*_3$ NIEOKREŚLONA; $G_3>R_{31}$ \textbf{wszędzie}; v38--v39 \\
\bottomrule
\end{tabular}
\end{center}

% -------------------------------------------------------------------
\subsection{Status i~spójność z~duchem teorii}
\label{app:L-status}
% -------------------------------------------------------------------

\begin{center}\footnotesize
\begin{tabular}{@{}p{5.0cm}lp{4.0cm}p{2.2cm}@{}}
\toprule
\textbf{Zdanie} & \textbf{Status} & \textbf{Warunki / założenia} & \textbf{Odsyłacz} \\
\midrule
Ogon solitonu przy zmiennym $\alpha_{\rm kin}$ jest oscylacyjny
  & \statuslabel{Twierdzenie}
  & linearyzacja ODE; $m_{\rm sp}^2 = V''(1) = 1$, niezależne od $\alpha_{\rm kin}$
  & stw.~\ref{prop:K-tail-form} \\[3pt]

Istnieje $z_0(\alpha_{\rm kin})$: pierwsze zero $B_{\rm tail}$
  & \statuslabel{Propozycja}
  & numeryczna (\texttt{ex58}, \texttt{ex62}); różne $\alpha_{\rm kin}$
  & uw.~\ref{rem:J-verif} \\[3pt]

$F(\alpha_{\rm kin})$ ma dokładnie dwa zera w~$[2{,}0;\, 3{,}5]$
  & \statuslabel{Propozycja}
  & numeryc.\ \texttt{ex72}; zakres $[2{,}0;\,3{,}5]$; tolerancja $10^{-8}$
  & wyn.~\ref{res:L-alpha-stars} \\[3pt]

Suma $\alpha^*_1 + \alpha^*_2 = 2\pi - 11/10$
  & \textbf{\statuslabel{Hipoteza} (obalona)}
  & \texttt{ex84} ($r_{\rm max}{=}120$): 9307~ppm;
    \texttt{ex87v2} ($r_{\rm max}{=}100$): 22336~ppm;
    \texttt{ex87v3}: błąd metodyczny (okno skalowane);
    wyniki niezm.\ oknem spójne z~odchyleniem $>0{,}5\%$
    (O-L1: obalenie potwierdzone, $\alpha^*_\infty$ w~toku)
  & hip.~\ref{hyp:L-sum-formula} \\[3pt]

Ogólność formuły sumy dla $\beta_{\rm pot} \neq 1$
  & \statuslabel{Program}
  & wymaga \texttt{ex73} z~pełnym skanem
  & O-L2 \\[3pt]

$\alpha^*_{1,2}$ to efektywne sprzężenia leptonu $e$/$\mu$
  & \statuslabel{Hipoteza}
  & interpretacja; wymaga unifikacji z~Ścieżką~9 (O-L3)
  & hip.~\ref{hyp:L-alpha-eff} \\
\bottomrule
\end{tabular}
\end{center}

\paragraph{Spójność z~duchem teorii.}
Cały formalizm Ścieżki~10 opiera się wyłącznie na:
\begin{itemize}[nosep]
  \item[$N0$-2] $\alpha = 2$ (podstawa skalowania kinetycznego;
    skanowane $\alpha_{\rm kin}$ odbiega od $\alpha_{\rm TGP}=2$ o~$20$--$40\%$),
  \item[$N0$-5] $\beta = \gamma$ $\Rightarrow$ $V(g) = g^3/3 - g^4/4$,
    $V'(g) = g^2(1-g)$,
  \item złotej proporcji $\varphi$ jako naturalnej skali między sektorami
    topologicznymi (hip.~\ref{hyp:J-golden-topo}).
\end{itemize}
Żaden nowy parametr nie został wprowadzony.
Wynik $r_{21}$ nie jest wstawiany jako wolny parametr wejściowy,
lecz \textbf{wyznacza} wartości $\alpha^*_{1,2}$: jest testem spójności teorii.

\paragraph{Skrypty numeryczne.}
\begin{enumerate}[nosep,label=\textbf{ex\arabic*.}]
  \setcounter{enumi}{61}
  \item \texttt{ex62\_phi\_fixed\_point.py}: poszukiwanie $\alpha^*$
    metodą $\varphi$-skalowania i~skan mapy $F(\alpha)$.
  \item \texttt{ex63\_scattering\_phase.py}: faza rozpraszania
    $\delta(\alpha) = \mathrm{atan2}(B, C)$; warunek Bohra--Sommerfelda.
  \item \texttt{ex64}--\texttt{ex70}: weryfikacja $\alpha^*$, masy
    solitonu przy $z_0$, porównania metod wyznaczania $A_{\rm tail}$.
  \item \texttt{ex71}--\texttt{ex72}: wyznaczanie $\alpha^*$ przez
    skan $\delta(\alpha)$ i~metodę brentq.
  \item \texttt{ex73\_sum\_formula.py}: weryfikacja ogólności
    $S = 2\pi - 1 - \beta_{\rm pot}/10$ dla różnych $\beta_{\rm pot}$.
  \item \texttt{ex74}--\texttt{ex79}: precyzyjna analiza bifurkacji
    $\beta_{\rm krit}$ i~dokładne wyznaczenie $\alpha^*_{1,2}$.
  \item \texttt{ex80\_algebraic\_analysis.py}: analiza algebraiczna
    $S$, $D = \alpha^*_2 - \alpha^*_1$, iloczyn~$P$.
  \item \texttt{ex81}--\texttt{ex84}: weryfikacja krzyżowa sum,
    dS/dR, precyzja $\alpha^*$ przez różne metody.
  \item \texttt{ex85\_agamma\_phi0\_precision.py}: hipoteza
    $a_\Gamma \cdot \Phi_0 \approx 1$ (Warstawa~III), odrębna linia badań
    (dokumentacja: sek.~\ref{ssec:param-classification}).
\end{enumerate}
